\SourceLeaf{443}%
meeting of two lines. It is therefore clear that lines that were shortest lines on one surface remain shortest lines after transfer. It follows readily that if any polygon formed by shortest lines corresponds, as situated on the first surface, to a figure of zeniths on the auxiliary sphere whose area is $=A$, while the same polygon, as transferred by development onto the second surface, corresponds to a figure of zeniths on the auxiliary sphere whose area becomes $=A'$, the equation
\[
A=A'
\]
will always hold. Although this proof originally assumes that the figures are bounded by shortest lines, it is readily seen to hold in general, whatever the boundary may be. Indeed, if the theorem is independent of the number of sides, then for any polygon some or all of whose sides are not shortest lines, nothing prevents us from conceiving of another polygon with infinitely many sides, each of which is a shortest line.

It is also clear that every figure retains its area when transferred by development.

\Needspace{6\baselineskip}
We are now concerned with $4$ figures:
\begin{enumerate}
\item[1)] any figure on the first surface,
\item[2)] the figure on the auxiliary sphere corresponding to the zeniths of the first figure,
\item[3)] the figure on the second surface formed by No. $1$ under development,
\item[4)] the figure on the auxiliary sphere corresponding to the zeniths of No. $3$.
\end{enumerate}
\noindent\hspace*{1.2em}From what we have proved, $2$ and $4$ therefore have equal area, as do $1$ and $3$. If these figures are now taken to be infinitesimal, the quotient obtained by dividing $2$ by $1$ is the measure of curvature at the location on the first curved surface, and likewise the quotient of $4$ by $3$ at the location on the second. Hence the important theorem that
\begin{quote}
»in the transfer of surfaces by development, the measure of curvature at each point remains unchanged«.
\end{quote}
\noindent The same therefore holds for the product of the greatest and least radii of curvature.

For the plane, the measure of curvature is evidently everywhere $=0$. From this
